\documentclass[../unravbook.tex]{subfiles}

\usepackage{parskip}
\setlength{\parskip}{10pt}
\setstretch{1.15}

\setenotez{
  list-heading = {\chapter*{#1}\markboth{#1}{#1}}
}

\begin{document}
\begin{spacing}{1.2}

\setcounter{chapter}{4}
\setcounter{endnote}{0}

\chapter{Israel, De-Racialized} %% {{{
    %% end header ----------------------------
    
    
In the framework of ideological arguments deployed against Israel reviewed in this book, the overarching narrative that is the foundation of accusations of apartheid and genocide must now be placed, in this concluding chapter, front and center: ``Race.'' With the notion of Israel's centering of its Jewish identity qualifying it to be considered an \enquote{ethno-racial supremacist state,} as established by its Jewish antagonists(such as Peter Beinart) as well as others, \footnote{Sanchez, D. (5/26/2015) \textit{Shapen in Dispossession - How the State of Israel was founded on the Plunder of Palestine}. Antiwar.com. \textsc{url:} \url{https://original.antiwar.com/dan_sanchez/2015/05/25/shapen-in-dispossession/} (``Thus the Zionist dream of Israel, from early on, was a dream ``of an ethno-racial supremacist state that could only be established through a massive colonialist campaign of ethnic cleansing and dispossession.'')} the subject of its inherent and increasingly grotesque forms of racism has become, over multiple decades, both widely disseminated and deeply explored. In effect, the topic of Race (or what may be more precisely described by the topic of \enquote{Racism,} as the concept of Race is now widely seen as a social construct that has \enquote{no biological basis} \parencite{Sesardic-2010}, has become the shining beacon that calls back to port all boats that detest the very idea of Israel's existence. The spirit of these attacks is reminiscent of a sentiment expressed by high-ranking French military officers who issued an open letter to the leaders of their government urging them to do more to defend national patriotism: ``Today, some speak of racialism, indigenism, and decolonial theories, but through these terms, these hateful and fanatical partisans want racial war. They despise our country, its traditions, its culture, and want to see it dissolve by tearing away its past and its history.''\footnote{Fabre-Bernadac, J-P (4/21/2021). \textit{``For a return of honor to our leaders'': 20 generals call on Macron to defend patriotism}. Valeurs Actuelles. \textsc{url: } \url{https://www.valeursactuelles.com/politique/pour-un-retour-de-lhonneur-de-nos-gouvernants-20-generaux-appellent-macron-a-defendre-le-patriotisme} (Translated from French.)} Sadly, the full brunt of these fears, as expressed by these signatories for France, have now come to fruition for Israel.

Since the early 2000s, race-related nomenclature has been firmly established through association with the crime against humanity, now formalized and defined in international law through the passage of the Rome Statute, known as Apartheid.\footnote{IMEU. (02/02/2022). \textit{How Israel Practices Apartheid.} IMEU; Institute for Middle East Understanding. \textsc{url:} \url{https://imeu.org/article/how-israel-practices-apartheid-imeu-policy-backgrounder}.} As this book has explored, this crime is now seen as a shoe that fits all possible sizes when placed on Israel's foot, and not only with respect to its ongoing military presence in the West Bank (or Judea and Samaria, or the ``Occupied Palestinian Territories,'' depending on the preferred source) to deal with security threats that are overlooked by the Palestinian Authority.\footnote{The term ``West Bank'' has recently been replaced by ``Judea and Samaria'' in legislation pertaining to what the Israeli government has long characterized as the \enquote{Disputed Territories}. See, The Knesset; State of Israel, February 12, 2025. \textsc{url:} \url{https://main.knesset.gov.il/en/news/pressreleases/pages/press12225q.aspx}. However, the phrase ``Occupied Palestinian Territories'' remains ``the official UN nomenclature for the West Bank, East Jerusalem, and Gaza Strip'' and is widely considered the \enquote{most consistent with international law and human rights discourse.} Occupied Palestinian Territories | Emor Institute. \url{https://emorinstitute.org/glossary/occupied-palestinian-territories/}.} As this book has also explored, these labels are in fact deeply inter-related, stemming from their coupled evolution into international law regimes as crimes against humanity deemed to be of the highest order,\footnote{Lingaas, C. (2015). \textit{The Crime against Humanity of Apartheid in a Post-Apartheid World}. Oslo Law Review, 2(2), pp. 86--115. p. 91. \textsc{url: } \url{https://doi.org/10.5617/oslaw2566} (``Although this ILC Draft coincides with the final years of the South African apartheid regime and as such certainly was influenced by the events of the time, it nevertheless demonstrates the uncontested seriousness of the crime of apartheid. Apartheid is repeatedly compared to genocide, which is nowadays considered 'the crime of crimes'.'')} But their conceptual sturdiness nonetheless depends on how reliable their shared perception of Palestinians, as members of a ``racialized'' population, ultimately is. This, however, must be considered an hypothesized abstraction as much as it might be perceived by those who embrace it as an unmistakable reality for the simple reason that, as noted in a prior chapter, ``defining precisely what is meant by racialization can be a daunting task, let alone identifying how, when, where, why, and under what conditions it operates.''\footnote{See Ch. 2, fn. 16.} 

Exploring the idea further also reveals that ``racialized'' ethnic groupings are differentiated from the ethnic group supposedly responsible for the racializing on the basis of power and privilege which, in the case of Palestine, devolves into those who speak Arabic as their first language (whether overwhelmingly Muslim or occasionally Christian in their confessionism). That is, differences in physical characteristics, often buttressed by accusations of an ``ethnic hierarchy among Israeli-Jews'' which favors ``Ashkenazi Jews ... from European countries or their offspring,''\footnote{Cf., Bernstein, D., Benjamin, O., \& Motzafi-Haller, P. (2011). \textit{Diversity in an Israeli intersectional analysis: The salience of employment arrangements and inter-personal relationships}. Women's Studies International Forum, 34(3), 220---231. p.220. \textsc{url: } \url{https://doi.org/10.1016/j.wsif.2011.01.011}. (``We use the term 'ethno-national' as the local parallel to race in the US context.'')} can only be seen as secondary, unless everyone below the purported apex of the ethnological hierarchy can also be considered, to some extent, racialized. This raises a broader issue, which is that as a theoretical construct that can be applied based on religion, language, gender or sexual preference, the widely acknowledged process of racialization muddles the meaning of race almost beyond comprehension because it no longer has any connection to what it was long based on, which was characteristics relating to human skin color and physiognomy.\footnote{Jablonski, N. G. (2020). \textit{Skin color and race}. American Journal of Physical Anthropology, 175(2), 437--447. \textsc{url: } \url{https://doi.org/10.1002/ajpa.24200}.} As noted recently by a Muslim scholar, this ``expansion of racism'' to include a multiplicity of ``structures and effects'' among ``intertwining threads in a fabric of oppression'' means that it has become a concept that is ``essentially postmodern in character.''\footnote{Caner K. Dagli (2/6/2020). \textit{Muslims Are Not a Race}. Renovatio | The Journal of Zaytuna College. \textsc{url: } \url{https://renovatio.zaytuna.edu/article/muslims-are-not-a-race}.} 

While postmodernism has been described as a ``highly contested [term] used to signify a critical distance from modernism.''\footnote{Buchanan, Ian (2018). \textit{Postmodernism}. A Dictionary of Critical Theory. Oxford University Press.} --- which may also be highly contested by those inclined to critically examine the worthiness of such broad philosophical categories --- there is at least a somewhat plausible view that postmodernists are ``suspicious of truth and objectivity and see claims to objective knowledge as a manifestation of a will to power or some other drive.''\footnote{Dagli 2020.} Their mode of thinking may be seen as having had its heyday in recent decades, particularly in the aftermath of World War II,\footnote{Kim, J. (10/11/2024) \textit{Postmodernism: An Idea Whose Time Has Come to Die}. Avant-Garde: A Journal of Peace, Democracy, and Science. \textsc{url: }\url{https://avantjournal.com/2024/10/11/postmodernism-an-idea-whose-time-has-come-to-die/}. (``Where did postmodernism come from? The simplest answer is the Second World War. Faced with the disaster of fascism and Nazism, European continental philosophy suffered a breakdown. A group of thinkers emerged who had little to no background in political movements, who had never struggled for anything—yet gained traction because they filled the ideological vacuum left by the wreckage of war and the failure of socialism in Western Europe.'')} but perhaps, like any other overarching epistemological viewpoint on how best to understand the world, it may now be susceptible to intellectual exhaustion on both sides of the divide between the notion that the values of Western civilization have now officially run their course and the view that they must be saved at all costs. In other words, and in the end, perhaps ``racializing'' Palestinians goes with the territory of vilifying Israel, but that does not mean that the territory is anything more a landscape of intellectual constructs. The making of essentially the same point can be detected in the introductory comments of a book on the ``politics of postmodernism'' from the 1980s (again, arguably in the midst of its generational heyday):\footnote{Ross, A. (ed.) (1986). \textit{Universal Abandon? The Politics of Postmodernism}. (Minneapolis: University of Minnesota Press) p. xi.}
 
\begin{adjustwidth}{.65cm}{.5cm}
    \begin{singlespace}
        [The fact that] the category of postmodernism  ... has achieved such diverse cultural currency as a term ... demonstrates what has been seen as one of [its] most provocative lessons; that terms are by no means guaranteed their meanings, and that these meanings can be appropriated and redefined for different purposes, different contexts, and, more important, different causes. In fact, this politics of appropriation, for so long exclusively the discursive preserve of the colonizer, has more recently been crucial to groups on the social margin, who have preferred, under certain circumstances, to struggle for recognition and legitimacy on established 'metropolitan' political ground rather than run the risk of ghettoization by insisting on the 'authenticity' of their respective group identities, ethnic, sexual, or otherwise.
    \end{singlespace}
\end{adjustwidth}

\begin{center}
    *\hspace{15pt} *\hspace{15pt} *
  \end{center}

This chapter yet proceeds on the tenuous assumption of an unalterable reality of a racial conflict between Israel and Palestine in both subtle and far-flung ways. This characterization has continued for decades, dating back at least to the passage by the UN General Assembly in 1975 of Resolution 3379, which announced that Zionism was a form of racism, so it is therefore worthy of a closer look as an element of the narrative of hatred toward Israel in its own right.\footnote{Elimination of all forms of racial discrimination: Zionism as racism - GA resolution. (03/11/2019). Question of Palestine. \textsc{url: } \url{https://www.un.org/unispal/document/auto-insert-181963/}.}  The ways in which racialization of the conflict surrounding Palestine has pervaded the deeply hostile climate now facing Israel can already be found peppered throughout this survey. They include heavy-handed as well as relatively oblique references, such as when a  UN summary on the nature and status of Indigenous peoples as one of many ``vulnerable groups'' around the world, covered within its website section on the adoption of the Universal Declaration of Human Rights, is presented under the categorical label of the need to ``Fight Racism.''\footnote{Chapt. 2, note 7.} From Nur Masalha's writings, we see the Zionist movement depicted repeatedly and in multiple ways as having been ``inspired by ... racial ideologies'' or inspired by ``ethnic'' or ``racial'' attitudes, which while certainly endemic across the world, are blandly described in Palestine's case as particularly ``problematic.''\footnote{Chapt. 2, note 6.} Masalha's relatively recent formulation of the key tenets of his thesis about Israel also echoes the conclusions of a Canadian-Palestinian citizen who proclaimed his belief in 2001 that ``Zionism is racism'' because it lacks an ``appreciation or acknowledgement of Palestinian ties to their own homeland,''\footnote{Chapt. 2, note 52.} which in turn mirrors UNGA 3379's declaration of Zionism as ``a form of racism and racial discrimination.''\footnote{Israel made the revocation of the UNGA's Resolution 3379 a condition for its participation in the Madrid Conference of 1991.United Nations General Assembly Resolution 46/86 (1991). Economic Cooperation Foundation: United Nations General Assembly Resolution 46/86 (1991). \textsc{url: } \url{https://ecf.org.il/issues/issue/1321}.} 

Also as noted earlier, Edward Said broadly reinforced this perception when he noted in in \textit{Orientalism}, that based on the crystallization of the Orientalists' ``system of truths'' in the nineteenth century, ``every European, in what he could say about the Orient, was consequently a racist.''\footnote{Chapt. 5, note 43.} In more recent years, this survey noted how a Canadian-Palestinian professor concluded that in assessing the overall relationship dynamics between Israelis and Palestinians, ``their racialized identities determine what they can know and how they can know it.''\footnote{Chapt. 7, note 2.} Toward the end of that same chapter, mention was made of an interview with Hamas leader Yahya Sinwar, in which he asserted that ``the same type of racism that killed George Floyd is being used by [Israel] against the Palestinians.''\footnote{Chapt. 7, note 55.} Clearly, then, when it comes to all the critiques of Israel by academics and Arab world public figures who seek its ultmate destruction, a key reason offered to justify that destruction is that while the rest of the world may have its problems with racism (as assuredly it does), Israel in particular must be considered spectacularly racist.  

In considering this widespread more carefully, the preceding chapter's review of the work of Patrick Wolfe may be the most informative in reflecting the direct connection between settler-colonialism and this widespread sentiment. The subject of racism may not seem to have been all that central to Wolfe's scholarship; that is, the primary focus of his writings is on the inevitability of indigenous erasure as ``a structure, not an event.'' But in his last full-length book in 2016,\footnote{Chapt. 8, note 69.} which he subtitled  \textit{Elementary Structures of Race,''} he concluded with an exhortations to his readership of to ``anti-racist collaboration'' and ``solidarities'' against it.\footnote{Chapt. 8, note 74.} This phraseology is instructive because it is here that we can begin to see what amounts to a closing argument of many of Palestine's most passionate advocates: that precisely because its predicament raises the specter of racism, it must be given tremendous --- indeed, global --- weight. To use one commentator's summary of the issue, ``Analyzing the workings of the color line in Palestine and Israel and mapping the trajectories through which race and racialization came to play a central role in the land, illuminate the truly global reach of its logic and highlight some of the properties that have made the global color line so durable.''\footnote{Bivins, A. (08/04/2021). Tracing the Historical Relevance of Race in Palestine and Israel. \textit{MERIP}. \textsc{url: } \url{https://merip.org/2021/08/tracing-the-historical-relevance-of-race-in-palestine-and-israel/}.} In other words, while it can be acknowledged that racism is everywhere, it is especially noteworthy in relation to Palestine and necessary to call out Israel for it.    

Clearly, then, the theme of Racism, with a capital ``R,'' is pivotal to this survey, epitomizing why Israel's adversaries have coalesced around it, in effect at the end of the ``color line'' rainbow, in their multitudinous protestations about indigeneity and settler-colonialism. It is all meant to lead to a shared, common and widely-held understanding that ``the Zionist entity has no right to `self-defence' against people under its occupation and siege''; that the ``Zionist project'' must be recognized as ``not only nationalism and settler colonialism,''; and that, the inescapable conclusion must be that it is ``firmly rooted in racial capitalism, and is a form of racism.''\footfullcite[][p. 242.]{Williams-2024} We thus perceive that notwithstanding its eventual revocation in 1991, UNGA 3379 has yet triumphed in having established the basis for the consensus view that global repudiation of Israel is a necessity. The same point was recently made by UN Special Rapporteur Francesca Albanese, when she noted that ``there is no war that Israel has ever waged against the Palestinians that can ever be deemed lawful. Israel has no right to wage a war to invoke self defense against the people it maintains under occupation.'\footnote{Open Source Intel (3/22/2025). Post. X.com. \textsc{url: } \url{https://x.com/Osint613/status/1903542856683307114}.}

This chapter will attempt to explore how claims pertaining to the inherent racist nature of Israel have become so firmly embedded into the indigeneity narrative around Palestine such that it results in the unalterable and absolute rejection of Zionism's legitimacy and, by extension, Israel's very right to exist. Identifying this as the existential issue --- the \textit{sine qua non} of the framing of the  Palestinian reality as one involving indigeneity in the first place --- reveals that it at least the penultimate driver (the ultimate one will be explored in the next chapter) toward the larger goal of building a diverse, powerful and global coalition of allies proudly engaged as a unified front against Israel. This coalition's obsessive goal is to oversee the overwhelming of the Jewish state with justified hatred resulting in violence, so that worldwide revulsion ultimately causes Israel's downfall as a nation-state. The precise manner of its demise is predicted to be similar to that of the Republic of South Africa before 1994, causing it to be transformed from a country that kept only its name (which Israel certainly would not be allowed to do) and was otherwise altogether re-engineered from the ground up as a result of external pressures. Investigating this as the end goal in the endless, exhaustive and extremist rationalizations contained within the vast corpus of literature on Palestinian indigeneity reflects the vast substantive breadth of Israel's predicament --- but for assembling all the pieces of the puzzle, as attempted in the next chapter and this survey's conclusion. It reflects why Angela Davis, an internationally renowned former Black Panther who has spent her entire adult life acting as an ``activist and educator within a politicized collective movement,''\footnote{Angela Davis (1944 ---) (historical profile). National Museum of African American History \& Culture. textsc{url} \url{https://nmaahc.si.edu/angela-davis}.} has repeatedly described Palestine as central to ``the development of solidarity campaigns that have brought different struggles together.''\footfullcite[][p. 42]{Davis-2016} 

Even so, because its implications are so far-flung --- that is, the intentional and premeditated destruction of an advanced, sophisticated and by all accounts largely prosperous nation-state with a population of over seven million citizens --- this desired end point is often acknowledged elliptically and without the finality of saying that Israel is irredeemable due to its inherent racism. That is, at least since the onset of the conflict resulting from October 7, 2023, it sees as thought the hand must be played without hitting people too much over the head with it. For example, mentions of issues pertaining to Race do not appear even once in a 2022 UN report by UN Special Rapporteur Francesca Albanese on the human rights situation in the Palestinian territories, even while the report goes to great lengths to eviscerate Israel as a settler-colonialist state.\footnote{United Nations General Assembly. (2022). Situation of human rights in the Palestinian territories occupied since 1967: Report of the Special Rapporteur (A/77/356). \textsc{url: } \url{https://www.ohchr.org/en/documents/country-reports/a77356-situation-human-rights-palestinian-territories-occupied-1967}} Also, while the word apartheid is nearly omnipresent in spoken or popular advocacy against Israel, such as in speeches, anecdotal writing and social media,\footnote{E.g., Robert, S. (8/12/2011). \textit{Apartheid on Steroids}. The Nation. \textsc{url: } \url{https://www.thenation.com/article/archive/apartheid-steroids/}.} analyses of the supposedly close correlation between Zionism and what may be characterized as ``Apartheidism'' --- that is, the presumed reality of any society deemed to be indistinguishable from the model of South Africa in its apartheid era --- are also conveyed in much drier fashion, without a strict definition of ``Race'' identified as a central issue.\footnote{Helou, Dania Ali (2018). \textit{Apartheid and the One-State Framework: Comparative Study of Israel and South Africa}. Masters Thesis, University of Memphis. \textsc{url: } \url{https://digitalcommons.memphis.edu/cgi/viewcontent.cgi?article=2976&context=etd}, p. 53 (citing Wolfe's 2006 article ``Settler colonialism and the elimination of the native,'' this first generation Palestinian-American master's degree candidate noted that ``the primary motivation of settler-colonialism for removing the indigenous population is not race, which is socially constructed and determined in its targeting, but the acquisition of territory.'').} In effect, the reasons to dwell on Palestine as the site of an Arabized yet also particularized indigenous people are to reach the conclusion that everyone worldwide should hold Israel accountable, and in contempt and  disgust, and thus be willing to eradicate it for the good of global order. Given the audience's familiarity with racism in the Western world as a topic of the utmost consideration, together with its supposed powerful antecedent slavery, all spawned by Europe's imperialist powers and its progeny (particularly the USA) of ``colorized, chattel slavery,''\footnote{Jackson, T. A. (2013). \textit{The great equalizer? Poverty, reproduction, and how schools structure inequality}. In P. C. Gorski \& J. Landsman (Eds.), \textit{The Poverty and Education Reader: A call for Equity in Many Voices} (pp. 117--131), p. 118. Routledge. \textsc{url: } \url{https://doi.org/10.4324/9781003448013}.} this is seen as the only plausible outcome.

But how has the concept of \textit{Race} become so seamlessly integrated and intertwined into the vast and growing corpus of modern scholarship on settler-colonialism? This involves complementary processes, neither of which are easily associated with Nur Masalha's emphasis on ``evidence-based history'' because they are largely derived from work done by the grounbreaking proponents of critical theory, as reviewed in the last chapter, together with what often seem to be deliberate misrepresentations of ``racial'' Zionist positions by Arab scholars.\footnote{Massad, J. (8/5/2019). \textit{How Zionists use racial myths to deny Palestinians the right to go home}. Middle East Eye. \textsc{url: } \url{https://www.middleeasteye.net/opinion/how-zionists-use-racial-claims-deny-palestinians-right-their-homeland}.(``Zionists do not claim that today's Egyptians, Syrians and Iraqis are pure descendants of the invading Arabs, rather than indigenous peoples who mixed with them. Yet, Zionists, like Netanyahu, insist fantastically that all Palestinians are newcomers to Palestine from the Arabian peninsula.'')} As such, their purpose is to ``critique the restrictive and alienating conditions of the status quo by analyzing the oppositions, conflicts and contradictions in contemporary society ... to eliminate the causes of alienation and domination (i.e., emancipate the oppressed class).''\footnote{Bhattacherjee, Anol, "Social Science Research: Principles, Methods, and Practices" (2012). Textbooks Collection. 3. \textsc{url: } \url{https://digitalcommons.usf.edu/oa_textbooks/3}, p.8.} Two recent articles by sociology and ``Africana Studies'' professor David Embrick,\footnote{Embrick, Dr. David G., Department of Sociology, University of Connecticut. \textsc{url: } \url{https://sociology.uconn.edu/person/david-embrick/}.}
who has focused throughout his career on the ``the impact of contemporary forms of racism on people of color,''\footnote{David G. Embrick, Loyola University Chicago, Department of Sociology. \textsc{url: } \url{https://www.luc.edu/sociology/faculty/Embrick,_D.shtml}.} illustrate the point. In both an earlier-cited and similar 2024 paper for Nur Masalha's \textit{Journal of Holy Land and Palestine Studies} and a 2023 paper along the same lines,\footfullcite[][]{Embrick-2023} racialization's role in how global social progress is made is met with stiff opposition by the forces of capitalism and colonialism, as best exemplified by none other than Israel:\footcite[][pp. 244-45, 253.]{Williams-2024} 

\begin{adjustwidth}{.65cm}{.5cm}
    \begin{singlespace}
        Israel, as an enduring settler-colonial enterprise, serves as a poignant microcosm that encapsulates many of the extremist and reactionary aspects of the global order — a legacy defined by centuries of colonialism, imperialism, exploitation, and conflict. Israel enforces a military occupation that sanctions, legitimises, and promotes attacks on and destruction of civilian populations, villages, and agricultural resources to acquire territory, strengthen its economy, and assert its religious dominance .... Examining Israeli apartheid and occupation through a racial capitalism lens provides a more comprehensive understanding of how Israeli racism, settler colonialism, and democracy-denial interlock with one another to dispossess and kill Palestinians. Without the racialization of Jews and Palestinians, the Israeli racial capitalism settler colonial state would not be possible.
    \end{singlespace}
\end{adjustwidth}

Intersectionality is an equally weighty concept, as underscored by the innumerable ``solidarity campaigns'' undertaken by those who see the commonality of reasons for struggle regardless of cause, so long as ``resistance'' and ``oppression'' can be applied to the bifurcated forces involved. As an academic concept, its contours were arguably first investigated --- or at least made subject to elaboration in a new way --- by UCLA professor Kimberle Crenshaw in her influential 1991 article \textit{``Mapping the Margins,''}\footnote{Crenshaw, K (1991). \textit{Mapping the Margins: Intersectionality, Identity Politics, and Violence against Women of Color}. Stanford Law Review, vol. 43, no. 6, pp. 1241--99. \textsc{url: } \url{https://doi.org/10.2307/1229039}} the struggles of ``women of color'' for being both women and of color are explored as being heavily intertwined. Recent authors have furthered intersectionality as parallel to the influence of Marxism in framing the political goals of the left.\footnote{Welchman, L., Zambelli, E., and Salih, R. (2020). Rethinking justice beyond human rights. Anti-colonialism and intersectionality in the politics of the Palestinian Youth Movement. \textit{Mediterranean Politics}, 26(3), 349--369. \textsc{url: } \url{https://doi.org/10.1080/13629395.2020.1749811}.} Thus, through a multiplicity of ``constant struggles of freedom,'' as envisioned by Marxist scholar Angela Davis, the higher goal of overthrowing interlocking systems of oppression is made clear, notwithstanding the apparent mediocrity of the idea itself that multiple oppressions (such as based on both sex and skin color) are incrementally worse than singular oppressions. In this discovery, the ``struggle for Palestine'' is now entrenched as a unidirectional imperative such that any sympathy for Israel's position vis-à-vis the ``resistance'' of Palestine is a non-starter. For those focused on this particular conflict, that may seem noteworthy, but the larger concern is found in the idea's universalist reasoning. It is all about struggle --- that is, the ultimate struggle --- against \textit{Weltschmerz}, meaning ``a feeling of sadness and lack of hope about the state of the world,''\footnote{Weltschmerz. Definition in the Cambridge English Dictionary. \textsc{url: } \url{ https://dictionary.cambridge.org/us/dictionary/english/weltschmerz}.} to use a German expression befitting the shared Marxist underpinnings of these ideas. 

Twenty years later, in a retrospective paper devoted to ``looking back to move forward,'' Professor Crenshaw noted that while progress has not yet been fully achieved in this intersectional, multifaceted struggle --- notwithstanding hopes for a post-racial utopia during the years of the Obama presidency --- it might yet still be on the horizon through the teaching of yet another variant in the outpouring of critical theory scholarship, Critical Race Theory, often referred to by the initials of ``CRT'':\footnote{Crenshaw, K. (2011). \textit{Twenty Years of Critical Race Theory: Looking Back To Move Forward}. Connecticut Law Review, vol. 43, no. 5, pp. 1253--1352, p. 1262.} 

\begin{adjustwidth}{.65cm}{.5cm}
    \begin{singlespace}
    The opportunity presented now is for scholars across the disciplines not only to reveal how disciplinary conventions themselves constitute racial power, but also to provide an inventory of the critical tools developed over time to weaken and potentially dismantle them. Beyond the academy, the opportunity to present a counter-narrative to the premature societal settlement that marches under the banner of post-racialism is ripe. In short, the next turn in CRT should be decidedly interdisciplinary, intersectional, and cross-institutional.
\end{singlespace}
\end{adjustwidth}

\noindent
A broad ``turn in CRT'' has indeed happened, with its theorists amply supported by ongoing efforts in the worlds of performing arts and popular media.\footnote{Examples are particularly plentiful in the work of popular ``rappers'' such as Vic Mensa's ``We Could be Free'' (\textit{url: } \url{https://youtu.be/SQqUuzIlb2k?si=FRC2kKwp-ilJRSDp}) and Benjamin Haggerty's ``fucked up'' (\textit{url: } \url{https://youtu.be/sn9EKC9nqU4?si=20WX1IXf-LJ8k2Sp}).} However, to proceed further either with an analysis of Palestinians and Israelis as ``racialized'' groups, whereby an effective intersectional strategy can be crafted in favor of one and against the other, definitional levels of meaning must be established to determine how groups can be identified as sufficiently ``racialized'' to establish that the bigotry toward them can be challenged by a countervailing strategy of \textit{``Anti-Racism,''} which is now frequently positioned as the only legitimate means by which the universalist goal of liberation from all forms of bigotry can be achieved.

\begin{center}
    *\hspace{15pt} *\hspace{15pt} *
  \end{center}

Thus far, this chapter has explored how ``Race'' and terms that are  etymologically derived from its broad understanding as the basis for structural systems of human oppression must derive from and be inter-related to a set of qualities that is stable enough to support a shared understanding of what the word and its offspring are intended to mean. Extended by racialization, the concept becomes very malleable and may now be found to encompass oppressions suffered in multitudinous ways by very different types of people, perceived along any number of vectors in relation to an oppositional type, which causes their diminution, dehumanization, or disadvantage. This includes people with traits that are physical, visible, tangible, emotional, identitarian, intellectual, or ideological, all so long as they can be considered deserving of a ``humanitarian'' necessity for respect of their differences from countervailing groups.\footnote{Hugo Slim (01/16/2018), \textit{Impartiality and Intersectionality}. Humanitarian Law \& Policy Blog. \textsc{url: } \url{https://blogs.icrc.org/law-and-policy/2018/01/16/impartiality-and-intersectionality/}.(``International humanitarian law (IHL) and humanitarian action know well that people's needs are structured by who they are, what has happened to them, and what power and vulnerability they possess. This means intersectional analysis of suffering and need, survival and recovery is intrinsic to humanitarian norms and practice.'')} If this respect is found wanting, then resistance may be all the more necessary to see that it is offered at a level and with a sufficiency deemed necessary by the person holding the differential identity. Described in such a way, the detection of racism, and the specification of who are the blatant and brutal racists and who are their victims, has become a finely tuned piece of machinery in ``international humanitarian law,'' so that whether the struggles are those faced by indigenous peoples or by ``people of color'' --- or by ``racialized'' peoples who fit neither category but can nonetheless be deemed oppressed by more powerful elements of society --- becomes the basis (at least by Embrick and scholars of similar ilk) for decisive counter-action. As thus described, they belong to what can be characterized as an expanded version of the ``Fourth World''\footnote{The ``Fourth World'' is a term that is widely attributed to the Tanzanian diplomat Mbuto Milando, who shared with George Manuel, a Canadian indigenous leader, that ``[w]hen native peoples come into their own, on the basis of their own cultures and traditions, that will be the fourth world.'' Manuel, G. (1974). \textit{The Fourth World}. Don Mills, Ontario: Collier-Macmillan Canada. p.5. However, it too has been broadened by some writers to cover all ``subpopulations existing in a First World country, but with the living standards of those in a third world, or developing country.''}  which in effect posits that all racism is a variant on the notion that all ``people of color'' are by their nature subjected to racism, regardless of whether ``color'' is the primary factor underlying their oppression. In this framework, ``[t]he life of an indigenous person in the Americas is rather similar to that of his black cousins, while holding different but also the same oppressions.''\footnote{Friesenhahn, C. (MerriCatherine), \& Kiksuya Khola (2018). \textit{An Introduction to the Fourth World: Does `working class' mean the same thing for all races?} The Anarchist Library. \textsc{url: } \url{https://theanarchistlibrary.org/library/merricatherine-and-kiksuya-khola-an-introduction-to-the-4th-world}. Kiksuya Khola appears to be a pseudonym of an anonymous person.} Leveraging not only Palestine's geographical proximity to the Western imagination given its connection to the ancient, biblical past, and because ``Jews are involved,''\footnote{Kant, L.  (7/24/2014). \textit{Gaza, Israel, and Media Coverage}. Mystic Scholar. \textsc{url: } \url{https://mysticscholar.org/gaza-israel-and-media-coverage/}. (``Israel/Gaza/West Bank is a tiny strip of land with a miniscule population. Even when there's no major conflict, the media focus is enormous and disproportionate. That's because it sells globally: in the U.S., in Europe, and in the Muslim world. It's because it's the land of the Bible. And it's because Jews are involved.''} Palestine is thus pushed to the top of the heirarchy of oppression, thereby making it emblematic of all others. 

At the same time, through the ``process of `being and becoming Palestinian' in a settler-colonial context [that] is the fundamental dynamic that molds present realities and generates future political possibilities,'' the necessity to ``de-exceptionalize the Palestinian condition by placing it in a global historical context, and to challenge the colonial genealogy of nationalist identity narratives by attending to the ironic and contingent historical pathways carved by discrete and internally differentiated Palestinian communities''\footfullcite[][pp. 9, 12.]{Doumani-2018} has grown to become a worldwide cause. Yet, notwithstanding this illusion of a world comprised of a monolith of oppressed peoples, all communities are in fact ``discrete and internally differentiated,'' but nonetheless collapsible as such by ``attaching meaning'' to [their] own group as a `race', and instilling within it  ``positive attributes,'' all as part of a tactic that has long been recognized as ``a common practice for subordinate groups seeking to defend and assert themselves collectively ... and instilling this meaning with positive attributes.''\footfullcite[][p. 20.]{Garner-2009} Even better, by asserting that Palestinians have been ``racialized'' is to establish that Zionism is ``racist,'' places their advocacy on straightforward ideological terrain, because the concept of ``race'' has become so malleable in contemporary discourse as the wellspring source of all human injustice, now far above the more narrow and now discredited set of economic problems relating to class, due to the suspicions many in the West tend to have toward classical Marxism as opposed to the modern ``Utopian'' basis for contemporary Identitarianism.\footnote{The distinction between classical Marxism, as described in Lenin's 1913 paper \textit{The Three Sources and Three Component Parts of Marxism}, and the critical theory that animates modern theorists who dwell upon the elusiveness of Utopia due to the dominance of Race, is explored in a recently published book by a Canadian professor of Philosophy who is Associate Editor for the journal \textit{Critical Philosophy of Race}. William Paris, \textit{Race, Time, and Utopia: Critical theory and the process of emancipation}  (1st ed). London: OUP. The book is described by its publisher as ``a theoretical account of utopia as the critical analysis of the sources of time domination and the struggle to create emancipatory forms of life.''}

Adapting this perspective, a broad consensus has emerged that racism is the result of an ``internalized white supremacy [which] affects white people and the dominant white culture in many ways,''\footnote{\textit{``Dismantling Racism: A Resource Book for Social Change Groups} (2023). Western States Center. \url{https://k-12-wa-public.s3.us-west-2.amazonaws.com/dismantling-racism-2003-western-states.pdf}, p. 49.} but with an end result that can be broadly described as ``the creation or maintenance of a racial hierarchy, supported through institutional power.''\footfullcite[][p.184.]{Kohli-2017}  Because ``racism'' and ``white supremacy'' are thus joined  as dichotomous, oppositional forces in a normative universe wherein there exists inextricable links between racism and the expansion of capitalism, colonialism and imperialism,\footfullcite[][p. 7.]{Batur-1999}  This results in a panoply of realities known as ``structural racism'' based on ``white supremacy,'' which becomes a layer of abstraction to denote ``the totality of ways in which societies foster racial discrimination through mutually reinforcing systems of housing, education, employment, earnings, benefits, credit, media, health care and criminal justice.''\footnote{What is Structural Racism? (9 Nov., 2021). The American Medical Association. \url{https://www.ama-assn.org/delivering-care/health-equity/what-structural-racism}.} The phenomenon of Racism takes on a meaning that is holistic of all other social science inquiry. Thus understanding its vastness, and accordingly dispensing with the need to place too much emphasis on exacting precision, a plausible enunciation of the core meaning of \textit{Racism} has been usefully offered by California professor Ramon Grosfoguel (partly by citing his own prior work) in 2016:\footfullcite[][p. 10.]{Grosfoguel-2016}

\begin{adjustwidth}{.65cm}{.5cm}
    \begin{singlespace}
Racism is a global hierarchy of superiority and inferiority along the line of the human that have been politically, culturally and economically produced and reproduced for centuries by the institutions of the ``capitalist/patriarchal western-centric/Christian-centric modern/colonial world-system.'' The people classified above the line of the human are recognized socially in their humanity as human beings and, thus, enjoy access to rights (human rights, civil rights, women rights and/or labor rights), material resources, and social recognition to their subjectivities, identities, epistemologies and spiritualities. The people below the line of the human are considered subhuman or non-human; that is, their humanity is questioned and, as such, negated (Fanon 1967). In the latter case, the extension of rights, material resources and the recognition of their subjectivities, identities, spiritualities and epistemologies are denied. (Citations omitted)
    \end{singlespace}
\end{adjustwidth}

But while this definition refers to the intellectual legacy of Frantz Fanon, it is perhaps even more grounded in the much earlier work of W.E.B. Du Bois, a Black educator and activist who was highly influential in the responses of Black Americans to the endemic racism of early twentieth century American life,\footfullcite[For a recent book describing how what CRT scholars describe as ``whiteness'' dominated the development of shared knowledge in American society until at least a century ago, \textit{see}][]{Lewis-2024} and who --- perhaps after recognizing what he considered to be the futility of opposing it by didactic means --- moved from a socialist frame of mind toward a full embrace of Marxism and Communism over the final decade of his long life.\footnote{From the People's World archives: W.E.B. DuBois joins the Communist Party. Feb. 23, 2021. \textsc{url: } \url{https://www.peoplesworld.org/article/from-the-peoples-world-archives-w-e-b-dubois-joins-the-communist-party/}.} But much earlier, in his 1903 ``groundbreaking book of racial philosophy'' entitled \textit{The Souls of Black Folk}, Du Bois established a view that has found footing in the exposition of multiple strands of critical theory --- and perhaps CRT in particular --- ever since: ``the problem of the Twentieth Century is the problem of the color-line.''\footnote{William Edward Burghardt Du Bois  and Manning Marable, \textit{The Souls of Black Folk} (originally published 1903; New York: Routledge, 2015). For the description of this book as ``groundbreaking,'' \textit{see} Jasmine Weber (02/05/2019). How W.E.B. Du Bois Meticulously Visualized 20th-Century Black America. \textsc{url:} \url{https://hyperallergic.com/476334/how-w-e-b-du-bois-meticulously-visualized-20th-century-black-america/}.} Du Bois elaborated on this assertion by noting that since the Middle Ages, ``the white races have had the hegemony of civilization,'' which they have used to perpetuate the ``color line'' just as robustly as they did centuries earlier.\footfullcite[][p. 6.]{Feagin-1999} 

Applying this approach to Israel as the ``white'' oppressor has even led some diasporic Jewish commentators to agree, contending --- much like Patrick Wolfe --- that the point is proven by the fact that \textit{Mizrachi} Jews who emigrated to Israel from Arab lands are generally more sympathetic to the plight of Palestinians than their Israeli \textit{Ahkenazi} neighbors because ``[their] goal was not to fight the idea of a Jewish homeland per se but to fight Ashkenazi Zionism, which they saw as intrinsically colonial and racist.''\footnote{Samuel, S. (04/11/2024). \textit{The untold story of Arab Jews — and their solidarity with Palestinians}. Vox. \textsc{url: } \url{https://www.vox.com/world-politics/24122304/israel-hamas-war-gaza-palestine-arab-jews-mizrahi-solidarity}. Considering herself a descendant of Jews from Arab lands, Ms. Samuel has also written about her pride in ``the progressive legacy of Mizrahi-Palestinian intersectionality'' because ``let's remind each other that intersectionality isn't some new, foreign implant into the Israeli-Palestinian conversation. It's only new if when you say `Israeli,' you really just mean `'mainstream, elite Ashkenazi.'"  Samuel, Sigal (03/13/2016). \textit{The Mizrahi-Palestinian Intersectionality Nobody's Talking About}. The Forward. \textsc{url: } \url{https://forward.com/opinion/335609/the-mizrahi-palestinian-intersectionality-nobodys-talking-about/}.} But to achieve this dichotomy of racist oppression, not only the oppressor must be identified --- that is, Israeli Jews who are the descendants of ``white'' European immigrants --- it is also necessary to identify holistically the oppressed population, in the form of all Arabs who identify themselves as Palestinians. Rather than consider Arabs in general or Palestinians in particular as a ``race'', claims of racism can thereby be levelled against anyone who supports the Jewish state (or at least anyone who is white within it or adapts a ``whiteness'' view, even if not white). If assessed more accurately in terms of the dynamics of their own intercultural relationships, they are more easily be characterized as people from a collection of local areas wherein Arab-speaking and acculturated clans have long been predominant (perhaps even ``oppressive'' power --- the Tawil and Barghouti in Ramallah, the Nashishibis and Husaynis in Jerusalem, the Yahya in Jenin, and various other places\footfullcite[][]{Cohen-2011} --- and other members of the Arab communities which they dominate are seen as their ``oppressed'' underlings. 

However, this view is overridden by virtue of their ``racialization,'' the basis for which may have originated in the work of Michael Omi and Howard Winant, who surveyed the phenomenon of ``racial formation'' in the United States, with their analysis holding true in all other societies that have disparate group with differing levels of economic, social and political power. In effect, the idea applies universally, although the United States provides a useful starting point for the template by which the same reasoning can be applied to other geographical areas such as Israel and Palestine: ``Race has been a fundamental axis of social organization ... that cannot be subsumed within a supposedly more fundamental category.''\footfullcite[][p.13]{Omi-Winant-1994} At the same time, Omi and Winant recognized that racial identity is not fixed but is instead ``an unstable and `decentered' complex of social meaning constantly being transformed by political struggle.''\footcite[][p.13]{Omi-Winant-1994} Indeed, the last decades of the twentieth century saw this tendency to define all oppressed groups engaging in political struggle racially play out nearly everywhere, so that they could be acknowledged as the ``resistance'' to the structural and cultural violence of colonialism, apartheid, and racial-ethnic cleansing. One result of these struggles is that in the words of Omi and Winant, ``we have now reached the point of fairly general agreement that race is not a biologically given but rather a socially constructed way of differentiating human beings,'' even while cautioning that ``the transcendence of biological conceptions of race does not provide any reprieve from the dilemmas of racial injustice and conflict nor from the controversies over the significance of race in the present.''\footfullcite[][pp. 55, 63.]{Omi-Winant-1994} In effect, the groundwork has thus been laid to engage in a process of ``racializing'' Palestinians because they can then be most clearly seen as the victims of the ``racial power'' of Israelis.  

\begin{center}
    *\hspace{15pt} *\hspace{15pt} *
  \end{center}

As racial disparities have remained a factor in the life of the United States and other advanced democracies of the western world, the emotional heft of the foregoing discussion on the importance of Race has expanded significantly. Over the last half century, this perhaps may be most clearly demonstrated by the career of Charles W. Mills (1951-2021), a political philosopher ``whose incisive criticism of liberalism and race both foreshadowed and framed contemporary debates about white supremacy and structural racism.''\footnote{Risen, C. (2021, September 27). Charles W. Mills, Philosopher of Race and Liberalism, Dies at 70. The New York Times. \textsc{url: } \url{https://www.nytimes.com/2021/09/27/us/charles-w-mills-dead.html}.}  As he began publishing books on racism in the 1970s, Dr. Mills sought to identify a direct line of descent from the 500-year era of imperialism and colonization to the modern era of structural racism. Possibly borrowing from the views expressed by Malcolm X in his speeches in the early 1960s,\footnote{X, Malcolm. (1989). \textit{The End of White World Supremacy: Four Speeches}. United Kingdom: Arcade Publishing} and drawing from his earlier attempts at ``theorizing white supremacy,''\footnote{Mills, C. W. (1994). Revisionist Ontologies: Theorizing White Supremacy. \textit{Social and Economic Studies}, 43(3), 105--134. \url{http://www.jstor.org/stable/27865977}.} he defined the problem bluntly in his 1997 book, \textit{The Racial Contract}: ``We live in a world which has been foundationally shaped for the past five hundred years by the realities of European domination and the gradual consolidation of global white supremacy.''\footfullcite[][p. 20.]{Mills-1997} In later years, Dr. Mills remained unimpressed by the growing literature on ``the facts of genocide and racism as narrative epicycles'' which he saw as amounting to little more than ``conceptual tokenization, where a black perspective is included but in a ghettoized way that makes no difference to the overall discursive logic of the discipline, or subsection of the discipline, in question.''\footnote{Lancaster, G. (2019, March 9) Review of the book \textit{``Black Rights / White Wrongs,''} by C.W. Mills. \url{https://marxandphilosophy.org.uk/reviews/16646_black-rights-white-wrongs-the-critique-of-racial-liberalism-by-charles-w-mills-reviewed-by-guy-lancaster/}.} Thus, since the reality of racism remains unaffected, such literature's publication was seen as irrelevant to the need to recompense those affected by the history's endless onslaught of past slights, all deduced to various forms and variants of racism.

\vspace{.1em}
Even further back than the era of Mills' writing career, the views of prominent German intellectuals from the middle decades of the last century serve as his outlook's antecedents. In Hannah Arendt's 1944 essay, \textit{``Race-Thinking before Racism,}''\footfullcite[][p. 36.]{Arendt-1944} she postulated that ``race-thinking'' was of recent vintage, attributable to Germans concluding they could not afford to be left out of the ``scramble'' by European nations to build far-flung empires. They thus justified their late entry into the global race for imperial conquests and colonial possessions by asserting their own racial superiority, entitling them to displace Slavs and destroy Jews who stood in their path. This racial-anthropological explanation of the hierarchy of importance among races soon became ``so powerful in so many fields of thought and practice''\footfullcite[][]{Larrimore-2006} that it now defined the intrinsic nature of Europe's colonial aspirations. This set the stage for today's scholars to contend that ``the colonial state exacerbates many racial and social relations within current Western societies'' such that ``each body is recognized as either colonizer or colonized.''\footnote{Books reviewed, author unlisted (2013). \# HISTORY /// The Desired Colonized Body: Foucault and Race by Ann Laura Stoler. \textsc{url: } \url{https://thefunambulist.net/editorials/history-the-desired-colonized-body-foucault-and-race-by-ann-laura-stoler}. In: The Funambulist.} 

We thus begin to detect the contours of the perceptual space necessary to understand the conceptual origins of how Palestinian cause is ultimately contextualized as an example of the horrors of \textit{Racism}, with a capital `R.' They are found in the need to challenge and ultimately defeat a quadruplet of evils gripping the world --- capitalism, imperialism, settler colonialism, and a resulting ``new racism'' that is structural in nature and attributable to ``white supremacy,''  through which all discernible winners and losers are inextricably linked to the hegemonic ambitions of the European countries that comprised ``the West''\footfullcite[This ``specific argument'' is made in][p. 88. (``{[}T{]}he new racism articulates ... the long history --- past and present --- of colonialism and racialization of those countries which participated and still participate in the production and reproduction of the `West' ...")]{Silva-1999} and their colonized but now independent offspring to subjugate other peoples in other lands and to maintain their post-imperial dominance, whether based on skin color or on other intangible attributes, such a shared legacy of migration from a Western European lands, even if the immigrants were themselves members of an oppressed population in such lands, as was the case for the vast majority of Ashkenazi Jews who made their way from Europe to Palestine in the twentieth century.\footnote{Anne-Maria Makhulu (2018). Book Review: Ann Laura Stoler, Duress: Imperial Durability in Our Times. In: \textit{Anthropological Quarterly}, Vol. 91, No. 4, p. 1451-1456, ISSN 0003-5491, p.1453 (``Are we to understand racism as subordinate to imperialism, as the consequence of imperial projects, as the condition of its making and durability? Or, perhaps, the elusiveness entailed in placing race and racism in relationship to empire is the intention; to complicate the determination between them and to assert that racism and empire continually invoke one another, thereby bringing one another into being.'')} In this narrative, those who are ``indigenous'' are reserved a special place among Black and Brown ``people of color,'' forming the ``I'' in the ``BIPOC'' acronym, because they possess the unique asset of ``indigenous knowledge.'' This is ``a rich social resource for any justice-related attempt to bring about social change'' as it relates to ``the dynamic way in which the residents of an area have come to understand themselves in relationship to their natural environment and how they organize their knowledge of flora and fauna, cultural beliefs, history, and teaching and learning to enhance their lives,'' so that ``indigenous ways of knowing become a central resource for the work of academics, whether they be professors in the universities or teachers in elementary and secondary schools.''\footfullcite[][p. 17.]{Kincheloe-2007} 

In the United States, Palestine's racialization in Western discourse has played an instrumental role within what has come to be recognized as the ``critical pedagogy'' movement, by which the tenets of critical theory are applied to turn the relative lack of opportunity of any ``racialized'' group vis-à-vis any other group as inferential proof of structural discrimination or racism, against which resistance must be directed toward the hegemonic nature of what is perceived to be the privileged, and therefore inherently oppressive, group. In the next chapter, the survey will take a closer look at how this worldview took root in education, causing Israeli responsibility for ongoing Palestinian suffering as a fundamental injustice that must finally be confronted so that Humanity can move away from the types of tribalistic loyalties that invariably stand in the way of an enlightened social liberalism that is necessary to achieve the fullness of its true potential. We can trace the arc of this narrative in the 2015 book by ethnic studies professor Keith Feldman, \textit{A Shadow Over Palestine: The Imperial Life of Race in America}, which makes the case that the racialization of Palestine became an integral part of American foreign policy over the 25-year period of 1960 to 1985, due to the centrality of Israel and Palestine in postwar U.S. imperial culture's identification of ``Islamofacism'' of all that was opposed to the values of ``liberal democracy'':\footfullcite[][]{Feldman-2015} 

\begin{adjustwidth}{.65cm}{.5cm}
    \begin{singlespace}
        Knowledge projects, cultural projects, and activist projects to elucidate the ineluctable humanity of Palestinians against their systematic exclusions developed relational analyses of racism, colonial violence, and imperial culture ... whose critical force registered desires for the enactment of substantive practices of decolonization in excess of the violent reproduction of U.S. and Israeli national exceptionalisms. In 1965, scholar-activists who were part of the PLO's Palestine Research Center worked to internationalize the Palestinian struggle by theorizing the particular and connected forms of racism animating Zionist settler colonialism. Black Power's Palestine envisioned Palestinian solidarity through the framework of connected anticolonial struggles for national liberation, touching down in places like the National Conference for New Politics, the Pan African Cultural Festival, and the United Front against Fascism Conference. 
\end{singlespace}
\end{adjustwidth}

Another paragraph on Feldman's deconstruction of Jewish neo-conservatism as a primary driver of this trend in official policy goes here.

In this narrative, a worldview that centers as an endpoint in contemporary social justice the paradigm of Palestinian suffering as ``racial'', framed within its indigenous nature and thus part of the intersectional ``BIPOC'' conglomeration which suffers ongoing and intractable racial oppression by the objective force of ''white supremacy,'' can be easily discerned, leaving as the only natural conclusion the goal of the oft-repeated chant: ``From the river to the sea, Palestine will be free.'' Opposition to such a vision of the future can thus only be described, according to a well-known Arab voice who envisions peace once world Jewry is stripped of its Zionist corruption (Yousef Munnayer, former executive director of the US Campaign for Palestinian Rights), as based on ``racism and Islamophobia.''\footnote{Yousef Munayyer (6/11/2021). What Does ``From the River to the Sea'' Really Mean? \textit{Jewish Currents}. \textsc{url: } \url{https://jewishcurrents.org/what-does-from-the-river-to-the-sea-really-mean}.} As will be further examined in the following chapter, the urgent need to finally ``defeat racism'' raises the stakes far beyond ``the river and the sea'' by many across the entirety of the Western world who are seen, according to their libertarian antagonists, as people convinced that they alone possess ``a kind of enlightened worldview adopted by good people.''\footnote{Deist, J. (2023). \textit{A Strange Liberty: Politics Drops Its Pretenses}. Ludwig von Mises Institute. \textsc{url: } \url{https://mises.org/online-book/strange-liberty-politics-drops-its-pretenses/15-getting-liberalism-wrong}. (``In this view, liberalism is a way of thinking, a kind of enlightened worldview adopted by good people. Liberalism becomes almost undefinable by its breadth: an interconnected passel of open-minded and tolerant attitudes across a swath of 'social' issues. Cosmopolitanism, reason, tolerance, feminism, antiracism, and a healthy degree of support for egalitarian democracy are at the fore; economics is not the focus. Yes, free markets and property rights play a supporting role, but modal liberalism is more concerned with human self-actualization made possible only with the correct kinds of left-cultural attitudes. Markets work, they make people better off materially, and most of all they yield the kinds of societies and people liberals like.''}

\parindent 0pt
\parskip 6pt

\end{spacing}

\setenotez{
  list-name = {Notes to Chapter \thechapter}
}
% Place this macro at the exact end of every chapter:
{\footnotesize
\printendnotes[paragraph]}

\thispagestyle{empty} % makes a blank page

\end{document}